\chapter{Exploración, árboles y caminos mínimos}
\label{md:exploracion}

\section{Matriz de adyacencia y recuento de caminos}
\label{md:adyacencia}

Ordenados los vértices $v_1,\ldots,v_p$, la \textbf{matriz de adyacencia}
de un grafo simple es $A=(a_{ij})$, con $a_{ij}=1$ si hay arista
entre $v_i,v_j$ y $0$ en otro caso. Es simétrica y tiene diagonal cero.
Para un digrafo simple, $a_{ij}=1$ si hay arco de $v_i$ a $v_j$;
ya no tiene por qué ser simétrica.
Renombrar y reordenar los vértices permuta simultáneamente filas
y columnas; una misma matriz representa grafos isomorfos.

Para $k\ge1$,
\[
 \boxed{(A^k)_{ij}=\text{número de caminos de longitud }k
          \text{ de }v_i\text{ a }v_j.}
\]
Aquí camino permite repetir vértices y aristas, según la
sección~\ref{md:caminos}. El producto matricial suma las posibilidades
para el vértice intermedio al concatenar caminos.

Si $i\ne j$, existe un camino de $v_i$ a $v_j$ si y solo si
\[
 (A+A^2+\cdots+A^{p-1})_{ij}>0.
\]
Basta buscar hasta longitud $p-1$, porque un camino puede simplificarse
sin repetir vértices. Para un grafo no dirigido con $p\ge2$, la
conectividad equivale a la positividad de todas esas entradas
fuera de la diagonal. Cada vértice está conectado consigo mismo
por el camino de longitud cero, aunque esa suma tenga una entrada
diagonal nula.

\begin{ejemplo}
Para las aristas $\{1,2\},\{2,3\}$,
\[
 A=\begin{pmatrix}0&1&0\\1&0&1\\0&1&0\end{pmatrix},\qquad
 A^2=\begin{pmatrix}1&0&1\\0&2&0\\1&0&1\end{pmatrix}.
\]
Hay un camino de longitud dos de $1$ a $3$, y dos de $2$ a $2$.
Los últimos repiten el vértice inicial y no son caminos simples.
\end{ejemplo}

\section{Árboles y árboles con raíz}
\label{md:arboles}

Un \textbf{árbol} es un grafo conexo sin ciclos.
Equivale a que haya exactamente un camino simple entre cada par
de vértices distintos: dos caminos distintos crearían un ciclo;
un ciclo permitiría elegir dos caminos entre algunos de sus vértices.

Todo grafo conexo finito contiene un \textbf{árbol generador}, un
subgrafo que es árbol y contiene todos sus vértices.
Se obtiene eliminando una arista de cada ciclo mientras quede alguno:
esa eliminación no rompe la conectividad.

Un \textbf{árbol con raíz} $(T,r)$ distingue un vértice $r$.
El camino único desde $r$ a cada vértice organiza los antecesores:
$v\ge w$ si $v$ está en el camino de $r$ a $w$, incluyendo sus extremos.
Es un orden parcial; dos ramas pueden contener vértices incomparables.
Orientar las aristas alejándose de la raíz expresa esa jerarquía.
El \textbf{nivel} de un vértice es la longitud del camino desde $r$;
un vértice no raíz de grado uno es una hoja.

\section{Algoritmo de Dijkstra}
\label{md:dijkstra}

En un digrafo finito con pesos \emph{estrictamente positivos}, el peso
de un camino es la suma de los pesos de sus arcos. Se busca el menor
peso desde una raíz $r$ a un destino, o a todos los vértices alcanzables.

Se asignan etiquetas tentativas $L(r)=0$ y $L(v)=\infty$ para $v\ne r$.
Inicialmente todos los vértices están pendientes. En cada paso:
\begin{enumerate}
 \item Elegir entre los pendientes un vértice $v$ con etiqueta mínima.
       Si su etiqueta es infinita, los restantes no son alcanzables.
 \item Fijar $L(v)$ como distancia definitiva. Si solo se busca un destino
       y $v$ es ese destino, se puede terminar.
 \item Para cada arco $v\to w$ con $w$ pendiente, actualizar
       \[
        L(w)\leftarrow\min\{L(w),L(v)+\text{peso}(v,w)\}.
       \]
       Si mejora, registrar $v$ como predecesor de $w$.
 \item Retirar $v$ de los pendientes y repetir.
\end{enumerate}
Los pesos positivos aseguran que pasar después por un vértice con
etiqueta mayor no puede mejorar la distancia ya fijada.
Los predecesores permiten reconstruir un camino mínimo retrocediendo
desde el destino hasta la raíz.

En grafos no dirigidos se actualizan los vecinos a través de cada
arista en el sentido del recorrido. Los empates pueden proporcionar
distintos caminos con el mismo peso mínimo.
Si todos los vértices son alcanzables, los predecesores forman un árbol
de caminos mínimos desde $r$: minimiza cada distancia desde la raíz,
lo que no equivale a minimizar la suma de pesos de todas las aristas
de un árbol generador.
